\scanpage{77}
\mspage{L}
% archive foliation: 76
pour les modèles normaux (\unsure{résultant} de
la partie facile de Lefschetz, du moins pour
les groupes fondamentaux rendus abéliens),
on a fini. De façon plus précise
il suffit, pour un \inserted{modèle} $X$ \uline{projectif} sur $k$,
(propre et normal), de prendre dans $X$
une ``courbe générique $C$'' et pour
$n(K)$ le genre de la dite courbe
générique.

\margin{\diagram{
\begin{tikzcd}
C' \arrow[r,phantom,"\subset"] \arrow[d] & X' \arrow[d] & \\
C_{k(\unsure{u|\omega})} \arrow[r,phantom,"\subset"] & X \arrow[r,phantom,"\subset"] & P^r
\end{tikzcd}
}}

Y-a-t-il une démonstration du fait
que cela donne une majoration, n'utilisant
pas la théorie du groupe fondamental
et le fait que $\uline{\mathrm{Pic}}^\tau_{X/k}$ est propre sur $k$
si $X$ est normale?
\margin{démonstration directe???}
